\uuid{jY2Y}
\titre{Équation homogène/solution particulière}
\auteur{hanani ; mzari}
\organisation{exo7}
\licence{CC BY-NC-SA 4.0}
\datecreate{2020-01-20}
\niveau{L1}
\module{Analyse}
\chapitre{Equation différentielle}
\sousChapitre{Résolution d'équation différentielle du deuxième ordre}
\difficulte{3}
\theme{Equations du second ordre, équation homogène, solution particulière}
\exo7qcmid{lille-2.equations-differentielles.027}

\contenu{
\begin{question}
\qclassification{225.02}
\qtags{motcle=équation homogène/solution particulière}
On considère les équations différentielles :
$$(E_1)\; :\; y''-y=3+\mathrm{e}^{2x}\quad \mbox{et}\quad (E_2)\; :\; y''-y=2+\mathrm{e}^{x}.$$
Parmi les affirmations suivantes, cocher celles qui sont vraies :
\begin{answers}
\good{$(E_1)$ admet une solution particulière sous la forme $\displaystyle y_0=a+b\mathrm{e}^{2x}$ avec $a,b\in \Rr$.}
\bad{$(E_2)$ admet une solution particulière sous la forme $\displaystyle y_0=a+b\mathrm{e}^{x}$ avec $a,b\in \Rr$.}
\bad{La solution générale de $(E_1)$ est : $\displaystyle y=-3+\mathrm{e}^{2x}+k_1\mathrm{e}^{x}+k_2\mathrm{e}^{-x}$, $k_1,k_2\in \Rr$.}
\good{La solution générale de $(E_2)$ est : $\displaystyle y=-2+\left(k_1+\frac{x}{2}\right)\mathrm{e}^{x}+k_2\mathrm{e}^{-x}$, $k_1,k_2\in \Rr$.}
\end{answers}
\begin{explanations}
Les solutions de l'équation caractéristique sont $\pm 1$. Donc $(E_1)$ admet une solution particulière sous la forme $\displaystyle y_0=a+b\mathrm{e}^{2x}$, car $2$ n'est pas une racine de l'équation caractéristique, et $(E_2)$ admet une solution particulière sous la forme $\displaystyle y_0=a+bx\mathrm{e}^{x}$ car $1$ est une racine simple de l'équation caractéristique.
\end{explanations}
\end{question}
}
